\textbf{COLUNI 2018 -- Questão 29}\par\medskip Em estudos na área da microbiologia, descobriu-se que a bactéria Streptococcus iniae afeta o cérebro de peixes, causando a “doença do peixe louco”. A partir de 1995, os criadores de trutas de Israel começaram a vacinar seus peixes. Apesar desta iniciativa, em 1997, ocorreu uma epidemia causada por uma linhagem de bactérias resistentes à vacina. Muitos cientistas acreditam que esta linhagem tenha surgido por pressão gradual seletiva induzida pela vacina.

\par

Este pressuposto sugere que a vacina permitiu:\begin{enumerate}
\item mutações nos peixes, tornando-os suscetíveis à infecção.
\item mutações nas bactérias, tornando-as resistentes à vacina.
\item aumento na proliferação da linhagem de bactérias mutantes resistentes à vacina.
\item desenvolvimento de anticorpos nos peixes, tornando-os resistentes à vacina.\end{enumerate}
